\documentclass[12pt]{article}
\usepackage{latexblog}

% ============================================================
% Blog Metadata
% ============================================================
\blogtitle{抓取QQ空间皮肤图片}
\blogdate{2017-01-11}
\blogtags{框架, Python, 闲话编程, Scrapy}
\bloglang{zh}
\blogsummary{}

\begin{document}
\makeblogtitle

最近把博客重新整理了一下，博文设置featured image果然看起来现代不少，但是要去哪找这么多合适的图片呢？当然PS是一个不错的选择，但是费时费力。看到QQ空间的皮肤倒是做的不错，直接拿来用吧，反正不违法。于是想用爬虫抓取。

先调试一下，找几个图片看看：

\begin{verbatim}
http://i.gtimg.cn/qzone/space_item/orig/3/103603_top.jpg
http://i.gtimg.cn/qzone/space_item/orig/7/101703_top.jpg
http://i.gtimg.cn/qzone/space_item/orig/5/108725_top.jpg
http://i.gtimg.cn/qzone/space_item/orig/15/111327_top.jpg
http://i.gtimg.cn/qzone/space_item/orig/8/106904_top.jpg
http://i.gtimg.cn/qzone/space_item/orig/10/102490_top.jpg
\end{verbatim}

可以看到前面的网址都是一样的\texttt{orig/\%d/\%d\_top.jpg}，刚开始以为前面的一个数字是主题编号还是什么的，后面自然是图片ID，于是用以下的脚本抓取（Scrapy):

\begin{Shaded}
\begin{Highlighting}[]
\KeywordTok{def}\NormalTok{ start\_requests(}\VariableTok{self}\NormalTok{):}
    \ControlFlowTok{for}\NormalTok{ i }\KeywordTok{in} \BuiltInTok{range}\NormalTok{(}\DecValTok{0}\NormalTok{, }\DecValTok{100}\NormalTok{):}
        \ControlFlowTok{for}\NormalTok{ j }\KeywordTok{in} \BuiltInTok{range}\NormalTok{(}\DecValTok{100000}\NormalTok{, }\DecValTok{119999}\NormalTok{):}
\NormalTok{            url }\OperatorTok{=} \StringTok{\textquotesingle{}http://i.gtimg.cn/qzone/space\_item/orig/}\SpecialCharTok{\%d}\StringTok{/}\SpecialCharTok{\%d}\StringTok{\_top.jpg\textquotesingle{}} \OperatorTok{\%}\NormalTok{ (i, j)        }
            \ControlFlowTok{yield}\NormalTok{ scrapy.Request(url}\OperatorTok{=}\NormalTok{url, callback}\OperatorTok{=}\VariableTok{self}\NormalTok{.parse)}
\end{Highlighting}
\end{Shaded}

后来抓取完成后，一共15个文件夹，每个文件夹差不多都是50-80个左右，这就有点意思了，可能腾讯利用了分布式的图片服务器，分散到0-15个不同的服务器上。随便找一个看看： 101195\_top.jpg在11下面，通常取余的方式实现，于是来试试：

\begin{verbatim}
101195% 16 = 11
\end{verbatim}

另外再验证几次也都是正确的，证明的我们的猜想。于是我们可以修改一下我们的爬虫了：

\begin{Shaded}
\begin{Highlighting}[]
\ControlFlowTok{for}\NormalTok{ i }\KeywordTok{in} \BuiltInTok{range}\NormalTok{(}\DecValTok{100000}\NormalTok{, }\DecValTok{200000}\NormalTok{):}
\NormalTok{    url }\OperatorTok{=} \StringTok{\textquotesingle{}http://i.gtimg.cn/qzone/space\_item/orig/}\SpecialCharTok{\%d}\StringTok{/}\SpecialCharTok{\%d}\StringTok{\_top.jpg\textquotesingle{}} \OperatorTok{\%}\NormalTok{ (i }\OperatorTok{\%} \DecValTok{16}\NormalTok{, i) }
    \ControlFlowTok{yield}\NormalTok{ scrapy.Request(url}\OperatorTok{=}\NormalTok{url, callback}\OperatorTok{=}\VariableTok{self}\NormalTok{.parse)}
\end{Highlighting}
\end{Shaded}

这样抓取就快多了，总共抓取了1232个图片。


\end{document}
